\begin{abstract}
Machine unlearning---the ability to remove specific data from trained models
without full retraining---is increasingly required by privacy regulation and
deployed in graph learning systems where deletion requests target individual
nodes or edges. Approximate graph unlearning (GU) methods offer dramatic
computational savings over retrain-from-scratch but introduce systematic
approximation errors. We ask whether these errors can be \emph{adversarially
exploited} by an attacker who controls deletion requests.

We present, to our knowledge, the first systematic adversarial audit of
approximate GU. The attacker selects a small set of nodes to forcibly unlearn
via influence-maximization, pseudo-influence functions, or hybrid strategies,
aiming to amplify approximation error and collapse downstream utility. We
evaluate five representative GU methods spanning three mechanism families
(learning-based, influence-function, shard-based) on Cora, Citeseer, and
ogbn-arxiv (169K nodes) with GCN and GAT backbones.

Our findings reveal a stark \textbf{vulnerability spectrum} determined by each
method's mechanism family.
(i)~Learning-based GNNDelete suffers F1 collapse averaging
\interim{16\%} with peaks above \interim{27\%} under a deletion budget
below 5\% of training nodes.
(ii)~Influence-function-based GIF shows a marginal but statistically
significant vulnerability (paired effect $\approx$\interim{+4\%} F1,
$p$\interim{$<$0.001}), bounded by its Newton-step approximation.
(iii)~Shard-based GraphEraser exhibits a counterintuitive \textbf{Shard
Protection Effect}---both random and adversarial deletion \emph{improve}
downstream F1 by \interim{3--6\%}, with the attacker's gain
upper-bounded by random selection, as the partition aggregator absorbs
deletion as regularization.
(iv)~Feature-preserving (MEGU) and certified-gradient (IDEA) families exhibit
null effect ($<$\interim{1\%} F1 difference), revealing a
\emph{mechanism-incomparable} robustness dimension that current selectors
cannot exploit.

We further introduce \textbf{collateral damage} diagnostics---retrain gap,
prediction shift, and a novel \textbf{hop-distance decay} metric quantifying
disturbance propagation through the GNN's receptive field---providing a
principled audit for any approximate unlearning deployment.
\end{abstract}
